\documentclass[../../../include/open-logic-section]{subfiles}

\begin{document}

\olfileid[hi]{sth}{cardinals}{zfc}	
\olsection{$\ZFC$: एक मील का पत्थर}

सुव्यवस्था जोड़ने के साथ हम अंतिम सैद्धांतिक मील के पत्थर पर पहुँच गए हैं।
अब हमारे पास $\ZFC$ के लिए आवश्यक सभी अभिगृहीत हैं। विस्तार से:

\begin{defn}
$\ZFC$ सिद्धांत के ये अभिगृहीत हैं: विस्तारिता, सम्मिलन, युग्म, घात
समुच्चय, अनंतता, आधारिता, सुव्यवस्था तथा पृथक्करण और प्रतिस्थापन योजनाओं के
सभी दृष्टांत। दूसरे शब्दों में, $\ZFC$, $\ZF$ में सुव्यवस्था जोड़ता है।
\end{defn}

$\ZFC$ का अर्थ \emph{चयन} सहित \emph{ज़रमेलो--फ़्रेंकल} समुच्चय सिद्धांत
है। यह कुछ विचित्र लग सकता है, क्योंकि हमने जो अभिगृहीत जोड़ा उसका नाम
``सुव्यवस्था'' था, ``चयन'' नहीं। किंतु बाद में जब हम {चयन} को सूत्रबद्ध
करेंगे, तो पता चलेगा कि सुव्यवस्था ($\ZF$ के सापेक्ष) चयन के तुल्य है
(\olref[choice][woproblem]{thmwochoice} देखें)। अतः इनमें से किसे अपना
``मूल'' अभिगृहीत मानें, इससे कोई अंतर नहीं पड़ता। और ``$\ZFC$'' नाम साहित्य
में पूर्णतः मानक है।

\end{document}
